\section{根sheafのEuler指数とcohomological degree}
\label{sec:sign-trichotomy}
Rainsの$\mathrm{ad}$は，一般に同じ非可換曲面の自己双対ではなく，
反対側の双対曲面$X^{\mathrm{ad}}$への反変導来同値である．
構造line bundle $e$はその双対$e^{\mathrm{ad}}=\mathrm{ad}(e)$へ移る．
既知の双対性
\[
 \mathbf R\mathrm{ad}_X:
 D^b(\mathcal A_{\mathrm R})^{\mathrm{op}}
 \xrightarrow{\sim}D^b(\mathcal A_{\mathrm R}^{\mathrm{ad}})
\]
とその二重双対を使う\cite[\S5, Corollary 5.9]{RainsBirational}．

\begin{lemma}[純一次元の双対と指数反転]
\label{lem:root-dual}
$M$を境界に制限して零となる既約な純一次元sheafとする．
$N=R^1\mathrm{ad}_X(M)$は双対曲面上の既約な純一次元sheafで，
\[
 \mathbf R\mathrm{ad}_X M\simeq N[-1],\quad
 R^1\mathrm{ad}_{X^{\mathrm{ad}}}N\simeq M,\quad
 c_1(N)=c_1(M),\quad \chi(e^{\mathrm{ad}},N)=-\chi(e,M).
\]
$M$上で自然な$\theta M\to M$が同型なら，
$N$上でも自然な$\theta^{\mathrm{ad}}N\to N$が同型である．
とくに，根sheafは双対でも根sheafになる．
\end{lemma}
\begin{proof}
Rains\cite[Corollaries 7.26, 7.30, 7.31]{RainsBirational}
によれば純一次元sheafでは$\mathrm{ad}$の導来cohomologyが
次数1だけにあり，その二重双対は元のsheafを復元する．
これにより$N$も純一次元である．
既約性はgenericな一次元support上の単純加群性に関する性質であり，
genericでの双対によって保たれる．
自然な境界triangleの双対は，双対側の境界triangleに一致するため，
自然変換が$M$上で可逆なら$N$上でも可逆である．
K群では$[\mathbf R\mathrm{ad}M]=-[N]$である．
Rains\cite[Proposition 7.16]{RainsBirational}の公式
$c_1(\mathbf R\mathrm{ad}M)=-c_1(M)$，
$\chi(\mathbf R\mathrm{ad}M)=\chi(M)$を使えば
$c_1(N)=c_1(M)$，$\chi(N)=-\chi(M)$となる．
ここで$\chi(N)=\chi(e^{\mathrm{ad}},N)$である．
\end{proof}

\begin{lemma}[零以下のEuler指数ではglobal sectionがない]
\label{lem:no-section}
$M$を根$\alpha^2=-2$を第一Chern類とする既約な純一次元sheafとする．
$m=\chi(e,M)\le0$なら$\Hom(e,M)=0$．
\end{lemma}
\begin{proof}
非零の$f:e\to M$があるとし，その像を$I$とする．
$M$の既約性は「任意の非零部分sheafが同じ第一Chern類を持つ」
ことであって，$f$が全射であることを意味しない
\cite[\S11, pp.~138--140]{RainsBirational}．
$I\ne0$だから$c_1(I)=\alpha$であり，商$T=M/I$は
零または零次元である．
$I$は$e$の商であるため，
Rains\cite[Corollary 10.16]{RainsBirational}の不等式
\[
 \chi(e,I)\ge-\tfrac12\alpha(\alpha+K_X)=1
\]
が適用できる．
一方，零次元sheafのEuler標数は非負
\cite[Proposition 7.6]{RainsBirational}なので
$\chi(e,M)=\chi(e,I)+\chi(e,T)\ge1$となり，
$m\le0$に矛盾する．
\end{proof}

\begin{theorem}[resonance sheafの符号別Extの完全判定]
\label{thm:sign}
$M$は$\alpha^2=-2$，$(-K_X)\alpha=0$を満たす根の
既約な純一次元sheafで，$Li^*M=0$とする．
$m=\chi(e,M)\in\mathbb Z$について全ての$n\in\mathbb Z$で
\[
\boxed{
\begin{aligned}
\dim\Hom(e,\tau^nM)&=\max(m,0),\\
\dim\Ext^1(e,\tau^nM)&=\max(-m,0),\\
\Ext^2(e,\tau^nM)&=0.
\end{aligned}}
\]
したがって自然な導来切断を$M$に適用すると，
\begin{equation}\label{eq:trichotomy}
 \Phi(M)=
 \begin{cases}
   \pi N_m,&m>0,\\
   0,&m=0,\\
   (\pi N_{-m})[-1],&m<0,
 \end{cases}
\end{equation}
ただし$m\ne0$のとき$\pi N_{|m|}\in\qgr R$は
境界$g=0$の外にある非零対象であり，
$R/(g-1)$の$|m|$次元右加群に対応する．
$m=0$のとき，$M$は非零であるが，全Serre orbitに
導来的に右直交し，指定した反標準導来切断は$M$を零に送る．
\end{theorem}
\begin{proof}
$Li^*M=0$なら境界triangleから自然な
$\theta M\to M$が同型であり，$\tau^nM\simeq M$．
line bundle $e$と一次元sheafに対する
$\Ext^2(e,M)=0$は
Rains\cite[Corollary 7.23]{RainsBirational}による．

$m\le0$では補題\ref{lem:no-section}により
$\Hom(e,M)=0$，Euler標数より$\dim\Ext^1(e,M)=-m$．
特に$m=0$では全$\Ext$が零となる．

$m>0$とする．双対$N=R^1\mathrm{ad}M$は
補題\ref{lem:root-dual}により双対曲面上の
負指数$-m<0$の既約根sheafで，境界にも制限して零である．
従って$\Hom_{X^{\mathrm{ad}}}(e^{\mathrm{ad}},N)=0$，
$\dim\Ext^1_{X^{\mathrm{ad}}}(e^{\mathrm{ad}},N)=m$．
Rainsの反変導来双対とSerre双対性を合成すると
\begin{equation}\label{eq:dual-ext}
\begin{aligned}
\Hom_X(e,M)
 &\cong\Ext^1_{X^{\mathrm{ad}}}(e^{\mathrm{ad}},N)^*,\\
\Ext^1_X(e,M)
 &\cong\Hom_{X^{\mathrm{ad}}}(e^{\mathrm{ad}},N)^*.
\end{aligned}
\end{equation}
実際$\mathbf R\mathrm{ad}_XM=N[-1]$なので
\[
 \Ext^{j}_{X^{\mathrm{ad}}}(e^{\mathrm{ad}},N)
 \cong\Ext_X^{j+1}(M,e)
 \cong\Ext_X^{1-j}(e,\theta M)^*
\quad (j=0,1).
\]
これに$\theta M\simeq M$を代入すれば
$\dim\Hom_X(e,M)=m$，$\Ext_X^1(e,M)=0$を得る．

最後に$g$は$\theta\to\mathrm{id}$から生じる中心正則元で，
$M$上で同型として作用するので，
各次数の$\Hom$または$\Ext^1$の$g$による移行は同型．
非零の次数付き$R$加群$N_{|m|}$は
$k[g]$上有限生成で，各次数の次元が$|m|$，
局所化の次数0部分は$|m|$次元の$R/(g-1)$右加群になる．
非零cohomologyが一つの次数に集中しているので
\eqref{eq:trichotomy}が従う．
\end{proof}

\begin{proposition}[root sheafは2-spherical]
\label{prop:spherical}
定理\ref{thm:sign}の$M$は任意の$m$について
\[
\Ext^j(M,M)=
\begin{cases}k&j=0,2,\\0&j\ne0,2.\end{cases}
\]
特に$\theta M\simeq M$の下で$2$-spherical objectである．
\end{proposition}
\begin{proof}
$M$の非零自己準同型の像は既約性により
第一Chern類$\alpha$を持つ．
核が非零なら零次元であり，$M$の純一次元性に反する．
よって自己準同型は単射であり，その余核は零次元．
自己射なので余核の数値類は零，Rains
\cite[Proposition 7.6]{RainsBirational}から余核も零．
従って$\End(M)$は有限次元の斜体で，$k=\mathbb C$だから
$\End(M)=k$．Serre双対により
$\Ext^2(M,M)=\End(M)^*=k$．
Rains\cite[Corollary 7.13]{RainsBirational}のEuler形式で
$\chi(M,M)=-\alpha^2=2$となるので$\Ext^1(M,M)=0$．
\end{proof}

\begin{corollary}[pointedなheart非一致と導来切断の限界]
\label{cor:pointed-sign}
$m<0$の既約根sheafがあると，Rains heartの$(e,\tau)$は
AZ-ampleでない．標準pair$(\pi R,(1))$がAZ-ampleなら
$(\mathcal A_{\mathrm R},e,\tau)$と
$(\qgr R,\pi R,(1))$はpointed abel圏として同値でない．
さらに指定した自然な$\Phi$は$m<0$で$t$-exactでなく，
$m=0$で保守的でない．
$m>0$の個々の根sheafの像はheartの次数0にあるが，
そのことだけでは$\Phi$の全体の$t$-exactnessは従わない．
\end{corollary}
\begin{proof}
$m<0$では非零$M$に対して全ての
$\Hom(e,\tau^nM)$が零なので，AZの生成公理が失敗する．
pointedなabel同値はample pairを保存する．
後二者は\eqref{eq:trichotomy}から従う．
\end{proof}
\begin{remark}[先行する証明の訂正と範囲]
「既約root sheafへの非零な構造sheafからの射は全射」
という推論は正しくない．必要なのは
像が同じ第一Chern類を持ち，商が零次元であることだけである．
補題\ref{lem:no-section}がその正しい証明を与える．
$m=0$の対象が実在することは共鳴の数値条件からは従わない．
genuine del Pezzo avatarでは後述の命題により$m=0$の根共鳴はない．
また指定した$\Phi$の非忠実性は，
別の抽象的導来同値の不存在を意味しない．
\end{remark}
